% Chapter 47: Spin, Angular Momentum Without a Classical Counterpart (Complete Chapter)
\chapter[Spin: Nonclassical Angular Momentum]{Spin: Angular Momentum Without a Classical Counterpart}

Previous chapters established amplitudes in Chapter~43, Schrödinger evolution in Chapter~44, state-changing measurement in Chapter~45, and vector/matrix accounts in Chapter~46. Now comes spin, quantum mechanics' famous and misunderstood angular momentum without anything rotating. Its two outcomes underpin magnets, hospital MRI, and quantum bits. By the end you should explain refrigerator magnets attracting nails and why spinning-electron pictures must fail.

\step{A Counterintuitive Opening}

% BEGIN TRANSLATOR SAFETY NOTE
\textbf{Translator safety note:} Strong magnets can pinch skin, shatter, and interfere with implanted medical devices. Keep them away from children and implants; prevent collisions and use eye protection. See \href{https://www.kjmagnetics.com/neodymium-magnet-safety.asp}{K\&J Magnetics' safety guidance}.
% END TRANSLATOR SAFETY NOTE

Buy an inexpensive neodymium magnet, the strong magnets found in refrigerator decorations, headphones, and phone speakers, and approach paper clips. It pulls the first through air, which lifts a second into an iron chain. Equal-sized copper or aluminum blocks do nothing however close.

Familiar since childhood, but why? Copper and iron both have atoms surrounded by negative electrons. Chapter~33 said moving charges make magnetic fields, and atomic electrons keep moving. Every atom should be magnetic; copper's electrons are no lazier. Where does magnetism hide?

% BEGIN TRANSLATOR SAFETY NOTE
\textbf{Translator safety note:} Treat the following cutting example only as a thought experiment. Do not saw, drill, or grind magnets: fragments and flammable dust pose hazards. See \href{https://www.kjmagnetics.com/faq.asp}{K\&J Magnetics' machining warnings}.
% END TRANSLATOR SAFETY NOTE

Saw a bar magnet in half and you obtain not separate north and south pieces but two complete magnets. Repeat, same result. Magnetism seems intrinsic to each electrically neutral atom. Perhaps each electron is a tiny magnetic needle? Correct. The next guess, a spinning charged sphere, causes a chapter of trouble. Magnetism leads to unfamiliar angular momentum that does not rotate, splits silver beams, and resides in your hydrogen nuclei awaiting MRI's roll call.

\step{Make a Guess First}

Write specific intuitions as usual.

First: picture an electron as a charged miniature Earth spinning about its center. Circulating charge makes current and magnetic moment, apparently explaining magnetism. What surface speed gives its measured moment: meters per second like a top, thousands like a projectile, or tens of thousands of kilometers near light speed? Choose a number for later calculation.

Second: evaporate silver into a beam through pointed and flat magnetic poles, creating a vertically varying field. Magnetic atoms are pulled differently according to north-up or south-up orientation. Does the screen show (a) a continuous vertical band from all possible angles, (b) two separate spots, or (c) one spot because atoms are too small to deflect? Record reasons.

Third: suppose two spots result. Select upper atoms and pass them through an \textbf{identically oriented} device. Do they split again? Rotate the next device ninety degrees for a horizontal gradient: what happens? Rotate back and remeasure vertically: does the initial selection survive? This cumbersome sequence is our sharpest knife.

\step{Hands-on Experiment}

Spin itself needs vacuum, a hot oven, and precise magnets, not home equipment. Instead examine the classical rotation we will leave behind.

\begin{experiment}[Swivel Chair and Compass: Classical Momentum and Magnets]
\textbf{Materials}: a smoothly rotating office chair or stool, two filled water bottles as weights, and a compass, including a phone app calibrated away from magnets.

% BEGIN TRANSLATOR SAFETY NOTE
\textbf{Translator safety note:} Do not perform the occupied-chair spinning and tilting steps below; use a recorded demonstration or the compass-only alternative. Chairs can tip or roll unexpectedly, and supporting someone does not remove that risk. See \href{https://ehso.emory.edu/sso/documents/toolbox-training_reducing-chair-related-injuries.pdf}{Emory Environmental Health and Safety's chair guidance}, which advises against leaning chairs or stools.
% END TRANSLATOR SAFETY NOTE

\textbf{Step one}: sit with arms extended and a bottle in each hand. A companion gently starts you rotating; lift your feet and pull both arms sharply inward. Rotation speeds up; extend and it slows. Angular momentum conservation compensates smaller inertia with greater angular speed, exactly a figure skater's technique.

\textbf{Step two}: notice angular momentum's axial direction. Carefully supported, try tilting yourself and the entire chair while rotating. Its axis resists. Direction matters as much as speed: which axis and which sense of rotation?

% BEGIN TRANSLATOR SAFETY NOTE
\textbf{Translator safety note:} The magnet precautions above apply to step three and the two-magnet alternative below: avoid snapping magnets together, keep them away from children and implanted devices, and protect your eyes. See \href{https://www.kjmagnetics.com/neodymium-magnet-safety.asp}{the manufacturer's safety guidance}.
% END TRANSLATOR SAFETY NOTE

\textbf{Step three}: place a compass on the table and record its direction. Approach with a neodymium magnet from various angles: the needle turns to align with the field; reverse the magnet and it reverses. Its moment seeks alignment with the external field. Crucially, direction is \textbf{continuously adjustable}: every field angle admits a matching needle angle, none skipped.

\textbf{Record and think}: classical angular momentum has continuously variable magnitude and direction, and magnetic needles can align stably in any direction. These are the old model's entire resources, both soon denied by microscopic angular momentum.
\end{experiment}

Without a chair, step three still helps. Approach two magnets and feel their two comfortable aligned, head-to-tail attracting orientations. Intermediate angles are possible if your hand maintains torque. Continuous orientation and torque define classical needles.

\step{The Old Model Runs into Trouble}

Calculate your spinning-electron guess.

Measured electron angular momentum is $\hbar/2$, with $\hbar\approx1.05\times10^{-34}$~J·s, roughly $5.3\times10^{-35}$~J·s. Classical rotational momentum is approximately mass times radius times surface speed, up to a shape factor near one. Electron mass is $9.1\times10^{-31}$~kg. Its radius? Colliders should reveal structure, yet deepest probes still see a point far smaller than $10^{-18}$~m. \textbf{Deliberately relax} this to classical radius $r\approx2.8\times10^{-15}$~m, the largest classical electromagnetic body available:
\[
v=\frac{L}{m\,r}\approx\frac{5.3\times10^{-35}}{9.1\times10^{-31}\times2.8\times10^{-15}}
\approx2\times10^{10}\ \text{m/s}.
\]
Light travels $3\times10^{8}$~m/s. Even this oversized electron needs \textbf{sixty or seventy times} light speed; actual radius bounds increase it thousands-fold. Chapter~38 forbids superluminal matter. Spin cannot come from bodily rotation, a fundamental failure starting sixtyfold beyond tolerance. Smaller radius only worsens it; a smaller electron cannot rescue the model.

Second, silver. In1922 Stern and Gerlach heated silver, collimated atoms through slits, sent them through pointed/flat magnetic poles, and deposited them on glass. Classical oven-emitted magnetic needles should have random orientations. Gradient force follows the field-direction moment component, continuously distributed, giving a vertical smear: fully aligned atoms at one extreme, anti-aligned at the other, all oblique cases between. This was almost everyone's expectation.

The result was \textbf{two} separated spots, upper and lower, with nothing between. No band or transition, as though only two directions existed. Did you predict that? Even Stern and Gerlach were surprised. Later understanding attributes almost all silver's moment to its lone outer electron, so beam splitting reveals electron magnetism and spin.

Third, the sequential questions. Block the lower beam and feed upper atoms into another identically oriented device: \textbf{all} go up, unsplit. A definite up property seems selected, like large beans passing the same sieve twice. Rotate the second ninety degrees: those up atoms \textbf{split again}, half left and half right. Still conceivable for needles: a vertical needle's horizontal component might have either sign. The shock comes after selecting left and sending it to a third vertical device. Classical ancestry certifies these as up atoms whose vertical property nobody changed, so all should go up. Instead, \textbf{half up, half down}. Horizontal measurement has \textbf{erased and rewritten} the vertical file, not read an existing property. Weighing selected large beans never turns half into small ones; sieving fails completely.

Chapter~45's polarizers already erased previous results this way: a forty-five-degree sheet recreated horizontal components from vertical photons. Silver repeats the same logic for angular momentum. This signals new state rules, not an apparatus quirk.

\step{Inventing New Concepts}

Redesign broken accounts one experimentally required step at a time.

First name the intrinsic angular momentum, fixed at $\hbar/2$ regardless of position, speed, or atomic binding, an electron identity like mass and charge: \textbf{spin}. The historical name suggests the very rotation just disproved. Larger composite-sphere rescues also failed. Chapter~49's Dirac theory makes spin \textbf{emerge automatically} when quantum mechanics meets relativity, with nothing rotating inside. Treat it as a proper name: as whales are not fish, spin is not spinning.

Second accept astonishingly few outcomes. Along any chosen direction, electron spin gives only up, $+\hbar/2$, or down, $-\hbar/2$. Vertical, horizontal, even forty-three degrees: always two. Classical moments give continuous components, for example total moment times $\cos30^\circ$, while spin is a two-sided coin tossable along any axis. Hence a \textbf{two-level system}, two state-space slots for one directional question, or \textbf{qubit}. Classical bits occupy 0 or 1; qubits superpose their basis states, a physical foundation for some quantum computers, revisited in Chapter~48.

Third convert states between directions. As polarization taught, each axis gives up/down basis states; expanding in another axis and squaring coefficients gives probabilities, then measurement updates the state. Initial vertical selection prepares up, equal left/right in the horizontal basis. Selecting left prepares a state equal up/down vertically, explaining the third split. Erased earlier files reflect two-level geometry, not rough instruments.

Fourth connect magnetism. Electron spin carries a moment about one Bohr magneton, $\mu_{\mathrm B}\approx9.3\times10^{-24}$~J/T, tied to spin direction, opposite for the negative electron, a sign detail not altering our conclusions. The gradient pulls aligned moments toward stronger field, pushes anti-aligned toward weaker, making two spots. Magnets, MRI signals, and magnetic storage trace back to this classically unmatched angular momentum.

Give the analogy, then rein it in. Spin resembles a coin's two outcomes with variable tossing axes. Unlike coins with preexisting orientation, electrons do not carry definite answers for every direction, Chapter~45's lesson. Chapter~48's Bell tests will defeat hidden prewritten answers mathematically too. Magnetic needle means only two projected moment results; actual needle moments come from macroscopic charge circulation with continuous directions, while electron moments are intrinsic and discrete in measurement.

\step{The Model in One Sentence}

\begin{oneline}
Spin is inborn angular momentum without rotation, fixed at $\hbar/2$ for electrons and yielding only along/against outcomes on any axis. Changing measurement direction invalidates old answers; its bound magnetic moment switches on magnetism and magnetic resonance everywhere.
\end{oneline}

Like birth registration, spin belongs to electron identity regardless of later experience. Unlike checkup measurements, $\hbar/2$ is identical for every electron; only question directions rotate. It is not miniature classical momentum. Your top's momentum is instead the macroscopic sum of immense numbers of spins and orbital motions.

\step{The Mathematical Translator}

Choose axis $z$ with states $|\,{\uparrow}_z\rangle$ and $|\,{\downarrow}_z\rangle$. An up state along a direction whose angle to $z$ is $\theta$ expands along $z$ as
\[
|\,{\uparrow}_\theta\rangle=\cos\frac{\theta}{2}\;|\,{\uparrow}_z\rangle+\sin\frac{\theta}{2}\;|\,{\downarrow}_z\rangle .
\]
Four questions: this expresses one state in another measurement basis. Coefficient squares sum to one and require pure up at $\theta=0$, pure down at $\theta=180^\circ$. The half-angle is spin's special rule, with origins beyond the main thread introduced in the challenge. It holds for ideal noninteracting single projective measurements; demagnetization, decoherence, and detector inefficiency make real apparatus approach that limit.

Born supplies our central number: atoms up along $\theta$ entering a $z$ Stern--Gerlach device emerge upward with
\[
P_{\uparrow}=\cos^{2}\frac{\theta}{2}.
\]
Check $\theta=0$: $\cos^2 0=1$, all up, repeated same-axis measurement unsplit. For $\theta=180^\circ$, $\cos^2 90^\circ=0$, all down, complete reversal. For $\theta=90^\circ$, $\cos^2 45^\circ=\tfrac12$, equal vertical outcomes from a horizontal state, matching the sequence. Notice the \textbf{half-angle}, versus polarization's full $\cos^2\theta$ in Chapter~45. This birthmark leads to spin states not returning unchanged after $360^\circ$.

For many identical $\theta$ atoms measured along $z$, the mean, in $\hbar/2$ units, is
\[
\langle S_z\rangle=\frac{\hbar}{2}\cos^{2}\frac{\theta}{2}-\frac{\hbar}{2}\sin^{2}\frac{\theta}{2}
=\frac{\hbar}{2}\cos\theta .
\]
At $\theta=0$, average $+\hbar/2$; at $\theta=180^\circ$, $-\hbar/2$; at $\theta=90^\circ$, zero. Individual outcomes remain $\pm\hbar/2$; continuous $\cos\theta$ belongs only to \textbf{repeated averages}. Discrete microscopically, continuous on average: the chair's adjustable angular momentum emerges this way.

\begin{figure}[htbp]
\centering
\begin{tikzpicture}[>=Stealth,scale=1.0]
  % Oven
  \draw[thick] (0,-0.55) rectangle (1.0,0.55);
  \node at (0.5,0) {Ag oven};
  \draw[thick,->] (1.0,0) -- (2.0,0);
  % Magnetic poles
  \draw[thick,fill=gray!25] (2.4,0.9) -- (3.6,0.9) -- (3.2,0.35) -- (2.8,0.35) -- cycle;
  \node[above] at (3.0,0.95) {N};
  \draw[thick,fill=gray!25] (2.4,-0.9) rectangle (3.6,-0.5);
  \node[below] at (3.0,-0.95) {S};
  % Split paths
  \draw[thick] (2.0,0) -- (2.6,0);
  \draw[thick] (3.6,0.15) .. controls (4.6,0.35) .. (5.8,0.55);
  \draw[thick] (3.6,-0.15) .. controls (4.6,-0.35) .. (5.8,-0.55);
  % Screen
  \draw[thick] (5.9,-1.0) -- (5.9,1.0);
  \fill (5.9,0.55) circle (0.06);
  \fill (5.9,-0.55) circle (0.06);
  \node[right] at (5.95,0) {Screen: two spots};
\end{tikzpicture}
\caption{Stern--Gerlach schematic. Pointed N and flat S poles create a gradient splitting silver by spin into two spots, not the classical continuous band.}
\end{figure}

Quantify the force. Let the vertical gradient be $10$~T per meter, readily achieved then. An atom with $\mu_{\mathrm B}$ feels approximately $F=\mu_{\mathrm B}\times(\text{gradient})\approx9.3\times10^{-24}\times10\approx9\times10^{-23}$~N. Silver mass $1.8\times10^{-25}$~kg gives $a=F/m\approx500$~m/s$^2$, fifty gravities. A few hundred meters per second crosses $0.1$~m in about $2\times10^{-4}$~s, deflecting $\tfrac12 a t^{2}\approx10^{-5}$~m, ten micrometers, with tens between beams. Visible glass spots turn ten-to-minus-twenty-three-newton forces on ten-to-minus-twenty-five-kilogram atoms into macroscopic evidence.

\begin{figure}[htbp]
\centering
\begin{tikzpicture}[>=Stealth,scale=1.0]
  \draw[thick,->] (-0.4,0) -- (10.6,0) node[right]{Atom beam};
  % Three devices
  \foreach \x/\lab in {1.6/{$z$ device}, 5.2/{Rotated $90^\circ$}, 8.8/{Back to $z$}}{
    \draw[thick,fill=blue!8] (\x-0.45,-0.8) rectangle ++(0.9,1.6);
    \node[below] at (\x,-0.85) {\lab};
  }
  \node[above] at (0.6,0.1) {Unselected};
  % Branches
  \draw[thick] (2.05,0.35) -- (3.1,0.6) node[above]{\shortstack{Up:\\keep}};
  \draw[thick,dashed] (2.05,-0.35) -- (3.1,-0.6) node[below]{\shortstack{Down:\\block}};
  \draw[thick] (3.35,0.6) -- (4.75,0.6);
  \draw[thick] (5.65,0.75) -- (6.7,1.05) node[above]{Left $\tfrac12$: keep};
  \draw[thick,dashed] (5.65,0.45) -- (6.7,0.15) node[below]{\shortstack{Right $\tfrac12$:\\block}};
  \draw[thick] (6.95,1.05) -- (8.35,1.05);
  \draw[thick] (9.25,1.2) -- (10.2,1.5) node[right]{Up $\tfrac12$};
  \draw[thick] (9.25,0.9) -- (10.2,0.6) node[right]{Down $\tfrac12$?!};
\end{tikzpicture}
\caption{Three successive Stern--Gerlach devices. Originally selected up atoms regain half down after horizontal measurement, which erases the earlier selection's file.}
\end{figure}

\begin{mathbridge}
\textbf{Pauli matrices: two-level algebra.} Write $|\,{\uparrow}_z\rangle$ and $|\,{\downarrow}_z\rangle$ as $\binom{1}{0}$ and $\binom{0}{1}$. Spin along three axes, in $\hbar/2$ units, corresponds to
\[
\sigma_x=\begin{pmatrix}0&1\\[2pt]1&0\end{pmatrix},\qquad
\sigma_y=\begin{pmatrix}0&-i\\[2pt]i&0\end{pmatrix},\qquad
\sigma_z=\begin{pmatrix}1&0\\[2pt]0&-1\end{pmatrix}.
\]
Build geometry: $\sigma_z$ has the basis eigenvectors and $\pm1$, up/down along $z$, hence diagonal. For $\sigma_x$, eigenvectors $\tfrac{1}{\sqrt2}\binom{1}{\pm1}$ equally superpose up/down, so $z$ states measured along $x$ split evenly. Conversely, $x$ eigenstates measured along $z$ also split evenly. Noncommutation $\sigma_x\sigma_z\neq\sigma_z\sigma_x$ has difference proportional to $\sigma_y$. Chapter~46 identified order-dependent accounts; file erasure is this matrix algebra. Generally,
\[
|\,{\uparrow}_{\theta,\varphi}\rangle=\cos\frac{\theta}{2}\,|\,{\uparrow}_z\rangle+e^{i\varphi}\sin\frac{\theta}{2}\,|\,{\downarrow}_z\rangle ,
\]
where spherical angles $\theta,\varphi$ locate a point on the Bloch sphere. Qubit states occupy its entire surface; classical bits have only the poles.
\end{mathbridge}

\begin{challenge}
\textbf{Two turns to return: spinors and $720^\circ$.} Half-angles imply a famous result. Advance the overall rotation angle, beyond parameter $\theta$, by $360^\circ$; $\cos(\theta/2)$ and $\sin(\theta/2)$ change sign, multiplying the vector by $-1$. Only $720^\circ$ returns it. One state's $-1$ leaves probabilities unchanged, an identity card rather than typo for \textbf{spinors}, half-angle-transforming objects. A macroscopic $360^\circ$ versus $720^\circ$ demonstration supports a book palm-up without overturning it while circling your arm around your body. One turn twists the arm; another untwists it. Topologically two-turn paths contract to rest, one-turn paths do not, Dirac's plate trick. Spin transformations doubly cover three-dimensional rotations. Deeper still, Dirac predicts magnetic-moment-to-spin $g$ factor $2$, Chapter~49; quantum electrodynamics corrects it to $2.002319\ldots$, agreeing beyond a dozen decimal places, among physics' most precise predictions. Its starting point is nonrotating angular momentum.
\end{challenge}

\step{The Model's Limits}

First, guard all rotating analogies. Electron-as-planet demands equatorial speed tens of times light, arithmetic suicide. Safely, spin is an intrinsic quantum number transforming under rotations by spinor rules, \textbf{contributing} angular momentum and obeying conservation without spatial circulation. A spin flip transfers the exact difference to photons or surroundings; accounts remain sound.

Second, up always needs an axis, like a mountain's height needing sea level or its foot. A definite $z$ state is a fully uncertain $x$ superposition, not incomplete knowledge. Perfect instruments still give $\pm\hbar/2$ with equal proportions.

Third, half-integer versus integer. Electrons, protons, and neutrons have $\hbar/2$; photons $\hbar$. This divides statistics: half-integer particles obey Pauli exclusion, forbidding two electrons the same state and explaining shells, periodicity, and solid floors. Integer particles can share states freely, as laser light exemplifies. Statistical chapters around Chapter~53 develop this. Here note that $\hbar/2$ ultimately determines why matter has volume, opacity, and resistance to collapse.

Fourth, MRI measures \textbf{ensembles}, not individual spins. Hospitals track no single proton; signals come from immense numbers of moments' \textbf{average magnetization}, precessing and relaxing. Two discrete individual outcomes yield continuous bulk waveforms. Reading each proton's state meets uncertainty; radio excitation and listening to bulk magnetization makes sense.

Fifth, spin has no consciousness or magical communication role. Flips transmit no information and individual outcomes are entirely random. Entangled spins in Chapter~48 still cannot signal superluminally. Quantum-resonance health scanners and spin-energy water have no connection to this physics.

\step{Explaining the Opening Phenomenon}

Why does neodymium attract the clips while copper does not?

All electrons have intrinsic moments; the difference is \textbf{pairing} and \textbf{alignment}. Most materials pair electrons in orbitals with Pauli-required opposite spins, canceling moments. All copper electrons are paired, making bulk copper nearly indifferent. Neodymium and iron have \textbf{unpaired} electrons; crucially, quantum exchange between neighbors, from Pauli and electrostatic repulsion rather than magnetic commands, favors \textbf{parallel spins}. Countless aligned spins each contribute $\mu_{\mathrm B}$, summing to a clip-lifting field. Iron's spins realign under neodymium's field, magnetization. Heating beyond the Curie temperature, about $770\,^\circ$C for iron, disrupts their ordered ranks, a thermal story.

Why complete magnets after sawing? Magnetism belongs to every spin, not the ends. Cutting divides one aligned army into two still-aligned armies. No magnetic monopole can be sawed out, and none has been experimentally found as of this writing.

% BEGIN TRANSLATOR SAFETY NOTE
\textbf{Translator safety note:} This remains a conceptual cutting example, not an instruction to cut a magnet. The \href{https://www.kjmagnetics.com/faq.asp}{manufacturer warns of machining and fire hazards}.
% END TRANSLATOR SAFETY NOTE

MRI is the same physics gloriously applied. Bodies are roughly sixty percent water, each molecule carrying two hydrogen nuclei, protons with $\hbar/2$ and moments. A typical $1.5$~T scanner field, thirty thousand Earth fields, gives along/against orientations with energy difference $\Delta E$ proportional to field. At $1.5$~T, proton precession is about $64$~MHz, radio frequency. Matching pulses push resonant protons like timed swing pushes, tipping bulk magnetization into precession and gradual recovery. Coils detect weak radio signals during recovery; tissue-dependent rhythms in fat, water, and tumors provide contrast. MRI calls a tuned group to sing, rather than photographs individuals. No rays or ionization, just two orientations and spin-bound moments. Turn off the field: $\Delta E\to0$, frequency vanishes, pulses select nothing, and imaging stops, expressing $f\propto B$ at $B=0$.

Earth's roughly $5\times10^{-5}$~T field does not show your compass quantum splitting. Its enormous ensemble and lattice produce continuous average magnetization, with thermal disturbances overwhelming individual gaps. Quantum rules are averaged, not absent: discrete microscopically, continuous macroscopically.

\step{Thought Experiments and Challenges}

\begin{thought}[Magnetars: Spin at Cosmic Extremes]
Neutron-star remnants pack more than a solar mass into roughly ten kilometers. Magnetars reach $10^{11}$~T, nearly a hundred billion MRI fields. Our ruler predicts proton precession about $4$~GHz, microwave, with electron splitting another three orders higher in soft X-rays. More visibly, fields confine atomic electrons into needles, rewriting bonds and deforming ordinary matter. The same $\hbar/2$ gives refrigerator magnets micronewton-meter torques and magnetars atom-distorting forces. Laws stay; parameters reach extremes.
\end{thought}

\begin{puzzles}
\item Someone gives electrons prewritten $x$, $y$, and $z$ answers, each $+$ or $-$, with measurement merely reading. Can sequential Stern--Gerlach results alone refute this? \textbf{Hint}: if the $z$ answer survives intervening $x$, repeating $z$ should retain up. What happens? This refutes only versions where $z$ is undisturbed by $x$ measurement; harder prewritten-answer theories await Chapter~48's Bell trial.
\item Select up from a $z$ silver beam, then use a device tilted from $z$ by $60^\circ$. What intensity ratio results? With a final $z$ device, using only initially selected up atoms, what follows? \textbf{Hint}: $P=\cos^2(\theta/2)$ gives for $\theta=60^\circ$ the $\cos^2 30^\circ=\tfrac34$ probability. Before the final answer, remember the middle measurement changed the state.
\item Can one electron spin serve as a continuously readable quantum gyroscope like a classical stable angular-momentum direction? Why use \textbf{many} spins, as in NMR gyroscopes? \textbf{Hint}: individual readings have only two outcomes. Where is continuous direction encoded? Recall discrete individuals and continuous averages.
\item Heat a neodymium magnet over an alcohol lamp and cool it; magnetism weakens. Do not keep strong magnets near flames too long; this is a thought question. Explain using spin and heat. Why no complete automatic recovery? \textbf{Hint}: heat randomly flips spins; cooling restores local order, but neighboring domains, independently aligned squads, need not agree.

% BEGIN TRANSLATOR SAFETY NOTE
\textbf{Translator safety note:} Do not heat a magnet or hold it over a flame, even briefly. Answer only theoretically: heating can ignite neodymium magnets. See \href{https://www.kjmagnetics.com/faq.asp}{K\&J Magnetics' heat and fire warnings}.
% END TRANSLATOR SAFETY NOTE
\end{puzzles}

Next place two spins together. Entanglement's spooky action faces Bell's experimental trial through these spins; quantum computing's magic is merely a many-partner dance on Bloch spheres.
